An overview of probability using measures. We will denote probability measures defined over $\sigma$-algebras with $\mathbb{P}$ and probability functions defined over some sample space $\Omega$ or $\mathbb{R}$ with $P$ or $p$. When we say countable, we mean finite or countably infinite. I have used resources from: 
\begin{enumerate}
  \item Rick Durret's \textit{Elementary Probability} and \textit{Probability} textbooks. 
  \item Dr. Krishna's \textit{Probability Foundation for Electrical Engineers} lectures at IIT. 
  \item Various quant interview books and websites for examples. 
\end{enumerate}


\section{Jan 12}  

  \begin{example}[Branching Processes]
    Let $Z_n$ be the number of people in the $n$th generation, and $X_i^{(n+1)}$ be the number of offspring the $i$th person in the $n$th generation has. Assume $X_i^{(n+1)}$ is iid with distribution $X$. Then, we have 
    \begin{equation}
      Z_{n+1} = X_1^{(n+1)} + X_2^{(n+1)} + \ldots + X_{Z_n}^{(n + 1)}, \quad Z_1 = 1
    \end{equation}
    What is the extinction probability $\pi$ for given $X$? 
    \begin{equation}
      \pi = \mathbb{P}(\exists n \text{ s.t. } Z_n = 0)
    \end{equation}
    Some ideas is that we can try to compute $\mathbb{E}[X]$, use a Markov chain. We can also use a generating function, which is defined 
    \begin{equation}
      f(\theta) = \mathbb{E} [\theta^X] = \sum_{k=0}^\infty \theta^k \mathbb{P}(X = k)
    \end{equation}
    Then, we can define $f_n (\theta) = \mathbb{E}[\theta^{Z_n}]$. Let $\pi_n = \mathbb{P}(Z_n = 0)$. By the dominated convergence theorem, 
    \begin{equation}
      \pi = \lim_{n \to \infty} \pi_n
    \end{equation} 
    We claim that $f_n (\theta) = f_{n-1} (f(\theta))$, which can be proven by the Tower property:
    \begin{align}
      f_n (\theta) & = \mathbb{E}[ \mathbb{E}[\theta^{Z_n} \mid Z_{n-1}]] \\ 
                   & = \mathbb{E}[ \mathbb{E}[\theta^{X_1^{(n)}} + \ldots + X_{Z_n - 1}^{(n)} \mid Z_{n-1}] ] \\ 
                   & = \mathbb{E}[f(\theta)^{Z_{n-1}}] \\ 
                   & = f_{n-1} (f(\theta))
    \end{align}
    where the final line is true since $f_{n-1} (\alpha) = \mathbb{E}[\alpha^{Z_{n-1}}]$. This gives us a nice recursive formula, which implies that 
    \begin{align}
      \pi_n = f_n (0) & = f_{n-1} (f (0))  \\
                      & = \underbrace{(f \circ \ldots \circ f)}_{n \text{ times}} (0) \\ 
                      & = f(f_{n-1} (0)) \\ 
                      & = f(\pi_{n-1}) 
    \end{align}
    So $\pi$ satisfies $\pi = f(\pi)$ as long as $f$ is continuous. 

    So let's actually do some computation. 
    \begin{enumerate}
      \item Assume $X \sim \mathrm{Geom}(p)$ i.e. $\mathbb{P}(X = k) = (1 - p)^k p$ for $0 < p < 1$. Then, 
        \begin{equation}
          f(\theta) = \frac{p}{1 - (1 - p)\theta} 
        \end{equation}
        Since $\theta = f (\theta)$, o we can solve 
        \begin{equation}
          \theta = \frac{p}{1 - (1 - p) \theta} \implies \theta = 1 \text{ or } \theta = \frac{p}{1 - p}
        \end{equation}
        So how do we know which solution is $p$? If $p \geq 1/2$, then $\pi = 1$, but if $\frac{p}{p - 1} < 1$, then we can derive further 
        \begin{equation}
          f_n (0) = \frac{p \big(\frac{1 - p}{p} \big) - p}{(1 - p) \big( \frac{1 - p}{p} \big)^n - p} \to \frac{p}{1 - p}
        \end{equation}
        as $n \to \infty$, so $\pi = \frac{p}{1 - p}$ is the correct extinction probability. 
    \end{enumerate}
    It turns out that according to some deeper theorem, if $\mathbb{E}[X] > 1$, then $\pi$ is the unique solution to $x = f(x)$ on $(0, 1)$, and if $\mathbb{E}[X] \leq 1$, then $\pi = 1$. 
  \end{example}

  \begin{example}
    We can describe $Z_n$ more precisely as $n \to \infty$ using a Markov chain, which can be stated in the equivalent ways. 
    \begin{enumerate}
      \item 
        \begin{equation}
          \mathbb{P}(Z_{n+1} = j \mid Z_0 = i_0, \ldots, Z_n = i_n) = \mathbb{P}(Z_{n+1} = j \mid Z_n = i_n)
        \end{equation} 

      \item 
        \begin{equation}
          \mathbb{E}[Z_{n+1} \mid Z_0, \ldots, Z_n] = \mathbb{E}[Z_{n+1} \mid Z_n] 
        \end{equation}
    \end{enumerate}

    Observe that $M_n = Z_n / \mu^n$ is a martingale, with $\mu = \mathbb{E}[X]$, i.e. $\mathbb{E}[M_n \mid M_{n-1}] = M_{n-1}$. So, 
    \begin{equation}
      \mathbb{E}[M_n \mid M_{n-1}] = \mathbb{E}[Z_n \mid Z_{n-1}] = \mu \frac{Z_{n-1}}{\mu_n} = \frac{Z_{n-1}}{\mu^{n-1}} = M_{n-1}
    \end{equation}
    By a Martingale convergence theorem, setting $M_\infty (\omega) = \lim_{n \to \infty} M_n (\omega)$, we have $M_n \to M_\infty$ a.s. 
  \end{example}

  Next time, we'll review conditional expectation, then martingales, then markov chains, then ergodic theory (building on convergence results), and finally with Brownian motion.

\section{Jan 14} 

	\begin{definition}[Conditional Expectation]
		Given probabiity space $(\Omega, \mathscr{F}_0, \mathbb{P})$, a $\sigma$-algebra $\mathscr{F} \subseteq \mathscr{F}_0$ and a random variable $X$ with $\mathbb{E}[X] < +\infty$, we define the \textbf{conditional expectation} of $X$ given $\mathscr{F}$---denoted $\mathbb{E}[X \mid \mathscr{F}]$---as \textit{any} random variable (called a \textbf{version}) $Y: \Omega \to \mathbb{R}$ s.t. 
		\begin{enumerate}
			\item $Y$ is $\mathscr{F}$-measurable. 
			\item For all $A \in \mathscr{F}$, $\int_A X \,d\mathbb{P} = \int_A Y \,d\mathbb{P}$
		\end{enumerate}
	\end{definition}
	\begin{proof}
		We should claim that it always exists and it is unique up to measure $0$. 
		\begin{enumerate}
			\item \textit{Uniqueness} Show that if $Y$ is a version, then $\mathbb{E}[|Y|] < +\infty$. Hint: Show that $\mathbb{E}[|Y|] < \mathbb{E}[|X|]$. We split $Y = Y+ - Y_{-}$. Then, $\mathbb{E}[|Y|] = \mathbb{E}[Y_+] + \mathbb{E}[Y_-]$, and we can bound 
				\begin{align}
					\mathbb{E}[Y_+] = \int_{\mathbb{R}_{>0}} Y_+ \,d\mathbb{P} = \int_{\mathbb{R}_{>0}} X \,d\mathbb{P} \leq \int_{\mathbb{R}_{\geq 0}} |X| \,d\mathbb{P} \\ 
					\mathbb{E}[Y_-] = \int_{\mathbb{R}_{<0}} Y_- \,d\mathbb{P} = \int_{\mathbb{R}_{<0}} X \,d\mathbb{P} \leq \int_{\mathbb{R}_{< 0}} |X| \,d\mathbb{P} 
				\end{align}
				Combining these together, we get $\mathbb{E}[|Y|] < \mathbb{E}[|X|]$. To prove uniqueness, we can see that if $Y, Y^\prime$ are two versions of $\mathbb{E}[X \mid \mathscr{F}]$, then let $A_\epsilon = \{\omega \in \Omega \colon Y(\omega) - Y^\prime (\omega) > \epsilon \}$. Then, 
				\begin{equation}
					\mathbb{P}(A) \epsilon \leq \int_A (Y - Y^\prime) \,d\mathbb{P} = \int_A (X - X) = 0 \implies Y \leq Y^\prime \text{ a.s.}
				\end{equation}
				Now swap $Y$ and $Y^\prime$ to get $Y^\prime \leq Y$. 
			\item \textit{Existence}. To prove existence of $Y$, this is much more involved. Recall the definition of absolutely continuous measures and $\sigma$-finite measures, along with the Radon-Nikodym theorem. We will use this to construct our conditional expectation. Since we can write $X = X_+ - X_-$, WLOG let $X \geq 0$. Let $\mu = \mathbb{P}$, we are looking for $Y$ s.t. for all $A \in \mathscr{F} \subset \mathscr{F}_0$, 
				\begin{equation}
				  \int_A Y \,d\mu = \int_A X \,d\mu
				\end{equation}
				Now set $\nu(A) \coloneqq \int_A X \,d\mu$, which turns out to be a $\sigma$-finite measure on $(\Omega, \mathscr{F})$. By construction, $\nu << \mu$. By the Radon-Nikodym theorem, $\exists f$ s.t. for all $A \in \mathscr{F}$, 
				\begin{equation}
				  \int_A f \,d\mu = \nu(A) = \int_A X \,d\mu
				\end{equation}
				Now set $f = Y$. 
		\end{enumerate}
	\end{proof}

	So it is the best approximation in a sense. The motivation is that it is a projection of $X$ onto the space of $\mathscr{F}$-measurable random variables. 

	\begin{example}
		Suppose $X, Z$ are random variables with a joint pdf $f_{X, Z} (x, z)$ and $\mathbb{E}[|X|] < +\infty$. Then, we define the conditional pdf as 
		\begin{equation}
			f_{X\mid Z} (x \mid z) \coloneqq \begin{cases} 
				\frac{f_{X, Z} (x, z)}{f_Z (z)} & \text{ if } f_Z (z) \neq 0 \\  
				0 & \text{ if } f_Z (z) = 0 
			\end{cases}
		\end{equation}
		Define $Y(z) \coloneqq \int_{\mathbb{R}} x f_{X \mid Z} (x \mid z) \,dx$. We claim that $Y(Z)$ is a version of $\mathbb{E}[X \mid \sigma(Z)]$. As a proof, for all $A \in \sigma(Z)$, we have 
		\begin{align}
			\int_A Y(z) f_Z (z) \,dz & = \int_A \bigg( \int_{\mathbb{R}} x f_{X \mid Z} (x \mid z) \,dx \bigg) f_Z (z) \,dz \\ 
															 & = \int_A \int_{\mathbb{R}} x f_{X, Z} (x, z)\,dx\,dz
		\end{align}
		where we denote $dx, dz$ are the Lebesgue or counting measure, depending on whether it is discrete or continuous. 
	\end{example}

	Now we list all of the properties of conditional expectation. 

	\begin{theorem}[Tower Rule]
		Let $X$ be a random variable on $(\Omega, \mathscr{F}, \mathbb{P})$ with $\mathbb{E}[|X|] < +\infty$. Let $\mathscr{G}, \mathscr{H}$ be sub-$\sigma$-algebras of $\mathscr{F}$. If $Y$ is any version of $\mathbb{E}[X \mid \mathscr{G}]$, then $\mathbb{E}[X] = \mathbb{E}[Y]$. 
	\end{theorem}
	\begin{proof}
	  $\Omega \in \mathscr{G}$ implies that 
		\begin{equation}
			\mathbb{E}[Y] = \int_\Omega Y \,d\mathbb{P} = \int_{\Omega} X \,d\mathbb{P} = \mathbb{E}[X]
		\end{equation}
	\end{proof}

	\begin{theorem}
		If $X$ is $\mathscr{G}$-measurable, then $\mathbb{E}[X \mid \mathscr{G}] = X$ a.s. 
	\end{theorem}
	\begin{proof}
		$X$ is a version of $\mathbb{E}[X \mid \mathscr{G}]$ because $X$ is $\mathscr{G}$-measurable. 
	\end{proof}

	\begin{theorem}[Linearity]
	  We have 
		\begin{equation}
			\mathbb{E}[a X_1 + b X_2 \mid \mathscr{G}] = a \mathbb{E}[X_1 \mid \mathscr{G}] + b \mathbb{E}[X_2 \mid \mathscr{G}]
		\end{equation}
	\end{theorem}
	\begin{proof}
	  We use linearity of the integral. For all $A \in \mathscr{G}$, 
		\begin{equation}
			\int_A (a X_1 + b X_2) \,d\mathbb{P} = a \int_A X_1 \,d\mathbb{P} + b \int_A X_2 \,d\mathbb{P} = a \int_A \mathbb{E}[X_1 \mid \mathscr{G}] \,d\mathbb{P} + b \int_A \mathbb{E}[X_2 \mid \mathscr{G}] \,d\mathbb{P}
		\end{equation}
	\end{proof}

	\begin{theorem}[Positivity]
		If $X \geq 0$, then $\mathbb{E}[X \mid \mathscr{G}] \geq 0$ a.s. 
	\end{theorem}
	\begin{proof}
		Let $Y$ be a version of $E[X \mid \mathscr{G}]$. If $\mathbb{P}(Y < 0) > 0$, then there exists a $\epsilon > 0$ s.t. $\mathbb{P}(Y < -\epsilon) > 0$. Set $A = \{\omega \in \Omega \mid Y(\omega) < -\epsilon\}$. Then, 
		\begin{equation}
			-\epsilon \mathbb{P}(A) > \int_A Y\,d\mathbb{P} = \int_A X \,d\mathbb{P} \geq 0
		\end{equation}
		a contradiction. 
	\end{proof}

	\begin{theorem}[Conditional Monotone Convergence Theorem]
	  If $0 \geq X_n$ is a monotonically increasing sequence of nonnegative functions and $X_n \to X$ a.s., then 
		\begin{equation}
			\mathbb{E}[X_n \mid \mathscr{G}] \nearrow \mathbb{E}[X \mid \mathscr{G}] \text{ a.s.}
		\end{equation}
	\end{theorem}
	\begin{proof}
		Let $Y_n$ be a version of $\mathbb{E}[X_n \mid \mathscr{G}]$. Set $Y(\omega) \coloneqq \limsup_n Y_n (\omega)$, which is $\mathscr{G}$-measurable and $Y_n \nearrow Y$. It remains to show that $Y$ is a version of $\mathbb{E}[X \mid \mathscr{G}]$. For all $A \in \mathscr{G}$, by MCT and the fact that $Y_n \geq 0$, we have 
		\begin{equation}
			\int_A Y_n \,d\mathbb{P} \to \int_A Y \,d\mathbb{P}
		\end{equation}
		But since $X_n \nearrow X$, $\int_A Y_n \,d\mathbb{P} = \int_A X_n \,d\mathbb{P} \to \int_A X \,d\mathbb{P}$, so $\int_A Y \, d\mathbb{P} = \int_A X \,d\mathbb{P}$, which implies that $Y$ is a version of $\mathbb{E}[X \mid \mathscr{G}]$. 
	\end{proof}

	\begin{theorem}[Conditional Fatou's Lemma]
	  If $X_n \geq 0$, then 
		\begin{equation}
			\mathbb{E}\big[ \liminf_n X_n \mid \mathscr{G} \big] \leq \liminf_n \mathbb{E}[X_n \mid \mathscr{G}] \text{ a.s.}
		\end{equation}
	\end{theorem}

	\begin{theorem}[Conditional Dominated Convergence Theorem]
		If $|X_n| \leq V$ for all $n \in \mathbb{N}$ and $\mathbb{E}[|V|] < +\infty$, and $X_n \to X$ a.s., then $\mathbb{E}[X_n \mid \mathscr{G}] \to \mathbb{E}[X \mid \mathscr{G}]$ a.s. 
	\end{theorem}

	\begin{theorem}[Conditional Jensen's Inequality]
		If $f: \mathbb{R} \to \mathbb{R}$ is convex, and $\mathbb{E}[|f(X)|] < +\infty$, then 
		\begin{equation}
			f\big( \mathbb{E}[X \mid \mathscr{G}] \big) \leq \mathbb{E}[f(X) \mid \mathscr{G}]
		\end{equation}
		As a corollary, we have $\|\mathbb{E}[X \mid \mathscr{G}]\|_p \leq \|X\|_p$ for $p \geq 1$, so conditional expectation is a contraction. 
	\end{theorem}
	\begin{proof}
		use the fact that $f(x) = \sup_n (a_n x + b_n)$ for some sequence $\{(a_n, b_n)\}_n$, and then invoke linearity. 
	\end{proof}

	\begin{theorem}[Tower Rule]
	  If $\mathscr{H}$ is a sub-$\sigma$-algebra of $\mathscr{G}$, then 
		\begin{equation}
			\mathbb{E}[\mathbb{E}[X \mid \mathscr{G}] \mid \mathscr{H}] = \mathbb{E}[X \mid \mathscr{H}]
		\end{equation}
		So, the smaller $\sigma$-algebra always wins. 
	\end{theorem}

	\begin{theorem}[Taking Out What's Known]
	  If $Z$ is $\mathscr{G}$-measurable and bounded, then 
		\begin{equation}
			\mathbb{E}[Z X \mid \mathscr{G}] = Z \mathbb{E}[X \mid \mathscr{G}]
		\end{equation}
	\end{theorem}
	\begin{proof}
	  Start out with indicators, then simple, then bounded as pointwise limits of simple functions. 
	\end{proof}

	\begin{theorem}[Role of Independence]
	  If $\mathscr{H}$ is independent of $\sigma(\sigma(X), \mathscr{G})$, then 
		\begin{equation}
			\mathbb{E}[X \mid \sigma(\mathscr{G}, \mathscr{H})] = \mathbb{E}[X \mid \mathscr{G}] \text{ a.s.}
		\end{equation}
		In particular, $\mathbb{E}[X \mid \mathscr{H}] = \mathbb{E}[X]$ a.s. 
	\end{theorem}
	\begin{proof}
	  Use $\pi$-systems and such, in Durrett's book. 
	\end{proof}


